% ECONOMICS_PROBLEM: {"id": "F51-430", "title": "Economic effectiveness of sanctions", "jel_code": "F51", "source_id": "OP-0086"}
% !TeX program = xelatex
\documentclass[11pt,a4paper]{article}
\usepackage[margin=23mm]{geometry}
\usepackage{fontspec}
\setmainfont{Latin Modern Roman}
\usepackage{amsmath,amssymb,mathtools}
\usepackage{xcolor,enumitem,booktabs,longtable,array,xurl,hyperref}
\definecolor{muted}{HTML}{53636C}
\hypersetup{colorlinks=true,urlcolor=blue,linkcolor=blue}
\setlength{\parindent}{0pt}
\setlength{\parskip}{5pt}
\newcommand{\R}{\mathbb R}
\newcommand{\E}{\mathbb E}
\newcommand{\N}{\mathbb N}
\newcommand{\ind}{\mathbf 1}
\newcommand{\Var}{\operatorname{Var}}
\newcommand{\Cov}{\operatorname{Cov}}
\newcommand{\supp}{\operatorname{supp}}
\DeclareMathOperator*{\argmin}{arg\,min}
\DeclareMathOperator*{\argmax}{arg\,max}
\newcommand{\entrymeta}[3]{{\sffamily\small\color{muted}\textbf{\texttt{#1}}\quad|\quad #2\par #3\par}}
\newcommand{\Problemref}[1]{\texttt{#1}}
\newcommand{\JELref}[2]{\texttt{#2} (JEL \texttt{#1})}
\begin{document}
\section*{Economic effectiveness of sanctions}
\textbf{Problem ID:} \texttt{F51-430}\quad \textbf{JEL:} \texttt{F51}\par
% BEGIN ECONOMICS STATEMENT
\label{problem:OP-0086}
\label{atlas:prob:T6}

\noindent\textbf{Atlas ID:} T6; \textbf{original number:} 86; \textbf{field:} International trade. \textbf{Classification:} Editorial research formulation (theoretical/empirical); not a certified unsolved theorem.

\subsection*{Definitions and assumptions}
Fix a finite-horizon game among a sanctioning coalition, a target government, firms and households. The coalition chooses a legally and technologically feasible sanction rule $s$; the target chooses actions $a_t$ including a prespecified political compliance criterion; firms choose production and rerouting within trade and financial restrictions. Model $\theta$ specifies payoffs, third-country responses, information, enforcement costs and exogenous states. Use a specified perfect Bayesian equilibrium selection with sequentially rational strategies and consistent beliefs. Let $C_H(s)$ be the indicator of compliance by horizon $H$; household utility and enforcement costs are separately measured and integrable.

\subsection*{Formal research question}
For regimes $s^0,s^1$ applied to the same initial state distribution, identify
\[
 \Delta_C=\mathbb E_\theta[C_H(s^1)-C_H(s^0)],
 \qquad
 \Delta_W=\sum_h\omega_h\mathbb E_\theta[U_h(s^1)-U_h(s^0)],
\]
where $\omega_h\ge0$ are explicitly chosen welfare weights. Characterize enforcement and network conditions that yield a positive compliance effect at a stated civilian and coalition cost. Develop credible partial-identification bounds if political compliance, evasion or the counterfactual no-sanction regime is incompletely observed. Determine whether sanctions that reduce direct bilateral trade remain effective once third-country diversion and expectations of future sanctions adjust.

\subsection*{Interpretation and scope}
Trade reduction, economic damage and political compliance are distinct outcomes. Compliance criteria must be set independently of realized success; retrospective redefinition creates selection bias. Treatment assignment is strategic: threatening states, target behavior and coalition commitment jointly affect observed sanctions. Identification therefore requires an exogenous assignment mechanism, defensible comparison restrictions, or a structural game with informative restrictions. Spillovers to third countries violate a simple untreated-control interpretation. Utility comparisons must specify whose welfare counts; a target-government payoff is not a measure of civilian welfare. Report sensitivity to equilibrium selection and the distribution of latent political shocks.

\subsection*{Source audit and status}
[S1] is a sanctions case-study book; its 2007 edition is verified by a PIIE author record, while product pages describe later formats. [S2] estimates trade effects of Russia-related sanctions; [S3] provides coded objectives and success. Neither trade effects nor coded success alone identifies $\Delta_C$. This is an editorial causal-policy question, not a certified open game-theoretic conjecture.

\subsection*{References}

\begin{enumerate}[label={[S\arabic*]},leftmargin=*]

\item Gary Clyde Hufbauer, Jeffrey J. Schott, Kimberly Ann Elliott, and Barbara Oegg (2007). \emph{Economic Sanctions Reconsidered}. 3rd edition, Peterson Institute for International Economics. \url{https://www.piie.com/bookstore/2016/economic-sanctions-reconsidered-3rd-edition-plus-cd-rom}. Date/version verification: \url{https://www.piie.com/experts/former-research-staff/barbara-oegg}.
\textit{Locator and role:} PIIE author bibliographic record (2007); publisher book contents, Chapters 2–4 (later format); PIIE author record verifies the 2007 third-edition citation; book contents concern sanctions case assessments. The available PIIE book-product pages describe 2008 and 2009 formats; they should not be mistaken for the 2007 edition date. Case success codings alone do not identify causal effects.

\item Matthieu Crozet and Julian Hinz (2020). \emph{Friendly Fire: The Trade Impact of the Russia Sanctions and Counter-Sanctions}. Economic Policy 35(101), 97–146. \url{https://doi.org/10.1093/epolic/eiaa006}.
\textit{Locator and role:} Publisher or author abstract / bibliographic record; Trade effects of the Russia sanctions and counter-sanctions; not general proof of political effectiveness. Provides background or a benchmark result; does not establish that the editorial question remains unresolved in 2026.

\item Gabriel Felbermayr, Aleksandra Kirilakha, Constantinos Syropoulos, Erdal Yalcin, and Yoto V. Yotov (2020). \emph{The Global Sanctions Data Base}. European Economic Review 129, article 103561. \url{https://doi.org/10.1016/j.euroecorev.2020.103561}.
\textit{Locator and role:} Publisher or author abstract / bibliographic record; Sanctions types, stated objectives and coded success during 1950–2016. Provides background or a benchmark result; does not establish that the editorial question remains unresolved in 2026.

\end{enumerate}

\subsection*{Common conventions: Source mathematical conventions}

Unless an entry states otherwise, agent and firm sets are finite. State and action spaces are measurable, strategies and outcome maps are measurable, probability laws are specified, and expectations exist for the targets under discussion. Infinite-horizon discount factors lie in $(0,1)$ and discounted objectives are finite. Norms, tolerances and complexity input sizes must be specified for computational comparisons. Constants or primitives that appear locally are inputs, not empirical facts asserted by this catalogue.

\textbf{Identification.} A statistical model $m\in\mathcal M$ implies an observable law $P_m$ and target $\theta(m)$. At law $P$, its identified set is
\[
\Theta(P;\mathcal M)=\{\theta(m):m\in\mathcal M,\ P_m=P\}.
\]
Point identification holds when this set is a singleton, or equivalently when $P_{m_1}=P_{m_2}$ implies $\theta(m_1)=\theta(m_2)$ for all compatible models. Sharp bounds mean bounds attained, or approached when the boundary is not attained, by compatible models. Writing this set is a definition, not a solution: the substantive task is to establish informative restrictions and derive or estimate the set. An unrestricted-confounding model may give only trivial bounds.

\textbf{Inference.} Identification, estimability and valid uncertainty quantification are separate questions. Fix $\alpha\in(0,1)$. A confidence procedure $C_n$ over a declared law class $\mathcal P$ has asymptotic uniform coverage at level $1-\alpha$ when
\[
\liminf_{n\to\infty}\inf_{P\in\mathcal P}\Pr_P\{\theta(P)\in C_n\}\geq1-\alpha.
\]
For set-identified targets the coverage event must be defined separately, for example coverage of the full identified set. If an entry uses identification-indexed models instead of uniquely indexed laws, the same coverage convention is applied over those models. Pointwise limits do not establish uniform validity, and asymptotic validity does not guarantee finite-sample precision. Sample sizes, dependence structures and weak-identification sequences are part of each application.

\textbf{Counterfactuals.} $Y_i(a)$ is a potential outcome under a specified intervention, including its timing, implementation and versions. It is not $Y_i$ conditional on voluntarily choosing $a$. A full intervention vector permits interference; shorthand scalar treatment notation does not silently impose no interference. Assignment, selection, attrition, anticipation and measurement must be specified. A claim about an instrument's local effect is not automatically a population policy effect.

\textbf{Economic equilibrium.} A counterfactual model includes best responses, constraints, feasibility, market clearing, state transitions and expectations consistent with the stated information. When several equilibria exist, either a declared selector is an input or the target is set-valued. Comparisons across policy regimes must use the stated selection convention. Reduced-form relationships are not presumed invariant after an equilibrium-changing intervention.

\textbf{Optimization and welfare.} Optimization is over the stated feasible set. If attainment is not established, $\sup$ or $\inf$ and $\varepsilon$-optimal choices replace an asserted maximizer or minimizer. For example, with value $V=\sup_{a\in\mathcal A}W(a)<\infty$, an $\varepsilon$-optimal set is $\{a\in\mathcal A:W(a)\geq V-\varepsilon\}$ for $\varepsilon>0$. Benefits, costs, risk and health or environmental outcomes can be aggregated only using declared units and conversion weights. Interpersonal utility weights, the treatment of fiscal transfers and foreign profits, discounting, and the behavioral welfare criterion are explicit inputs. A comparative-static derivative additionally requires existence and appropriate differentiability of the selected counterfactual.

\textbf{Sources.} Local labels [S1], [S2], and so forth restart in every question. Locators identify the relevant abstract, model, section or result at the level actually checked; a bibliographic or abstract-level verification is not represented as inspection of an entire proof. Working papers are distinguished from later publications and changing versions. Source summaries are paraphrased. All URL checks and editorial formulations refer to the audit date stated above.
% END ECONOMICS STATEMENT
\end{document}
